% BORRADOR (comentarios del director, 19-oct-2026): reemplaza MainMatter/Cap2/M03-Chapter2.tex. Borrar esta línea al pasarlo.
Este capítulo revisa la literatura en torno a los componentes de un pipeline RAG para la búsqueda de empleo. Empieza por el problema, la búsqueda léxica y sus límites, y por las soluciones semánticas de la industria del reclutamiento. Sigue con el fundamento de RAG y con cada componente que este trabajo evalúa: filtros estructurados, búsqueda híbrida, reordenamiento y modelos de \textit{embeddings} multilingües. Después revisa los motores de almacenamiento vectorial, la clasificación de consultas y la evaluación de sistemas RAG, y cierra con el marco metodológico. Cada sección termina con su relación con JobAI.

\section{Búsqueda léxica y sus límites en el reclutamiento}
\label{sec:estado-arte-keyword}

La búsqueda léxica relaciona una consulta con un documento por los términos que comparten. Su forma más simple exige que el texto de la consulta aparezca literalmente en el documento; su forma de referencia, Okapi BM25 \citep{Robertson2009}, pondera cada término por su rareza en la colección y normaliza por la longitud del documento. Ambas comparten el \emph{desajuste de vocabulario}: un sistema léxico no puede anticipar todas las formas en que dos personas describen el mismo concepto.

En el reclutamiento, el problema tiene dos caras. \citet{HaThuc2017} lo describen desde el reclutador: traducir criterios de contratación a una consulta booleana exige pericia y, aun así, omite candidatos que describen la misma experiencia con otras palabras. \citet{Bocharova2023} lo describen desde el catálogo de vacantes: dos títulos de cargo pueden ser funcionalmente equivalentes sin compartir ningún término, y un sistema de coincidencia exacta los trata como distintos. El resultado son dos errores que quien busca no distingue: ofertas relevantes que no aparecen porque no comparten el vocabulario de la consulta, y ofertas irrelevantes que aparecen por una palabra clave aislada.

Que la búsqueda léxica tenga esta limitación no la convierte en una línea base débil. \citet{Thakur2021} compararon diez arquitecturas de recuperación sobre 18 colecciones de dominios distintos y encontraron que BM25 es una referencia robusta: varios modelos densos no la superaron fuera de su dominio de entrenamiento, y los modelos de reordenamiento obtuvieron los mejores resultados a un alto costo computacional. Una evaluación que compare un sistema semántico solo contra una búsqueda por subcadena sobrestima su ventaja.

\textbf{Relación con JobAI.} La plataforma ofrece en su panel una búsqueda por subcadena sobre el título y la empresa, que la evaluación conserva como línea base de la propia plataforma. La comparación principal, sin embargo, se hace contra BM25, la línea base que la literatura considera fuerte (Capítulo~\ref{cap:metodologia}).

\section{Recuperación semántica en el emparejamiento de empleo}
\label{sec:estado-arte-empleo}

La industria del reclutamiento en línea adoptó pronto la búsqueda semántica con representaciones densas. \citet{HaThuc2017} documentan el paso de LinkedIn Talent Search de una búsqueda por palabras clave a una búsqueda por ejemplos: el reclutador indica uno o varios candidatos ideales y el sistema aprende a construir la consulta y el \textit{ranking} a partir de ellos. \citet{Bocharova2023} proponen VacancySBERT, una red siamesa basada en \textit{Sentence-BERT} que representa títulos de vacantes y habilidades en un espacio donde la similitud coseno refleja equivalencia funcional; reportan mejoras de hasta 21,5\,\% sobre modelos genéricos al incorporar las habilidades de la oferta, una evidencia de que una representación especializada en el dominio supera a una de propósito general.

A mayor escala, \citet{Shen2024} describen cómo LinkedIn incorporó la recuperación aprendida, densa y generativa, a su sistema de búsqueda y recomendación de empleo, en reemplazo parcial de un esquema de índice invertido por entidades extraídas. Más recientemente, la misma plataforma usa un modelo de lenguaje para interpretar la consulta del usuario y extraer de ella atributos estructurados que alimentan la búsqueda, en lugar de varios modelos de reconocimiento de entidades especializados \citep{Liu2025}. Otros trabajos usan un modelo de lenguaje que genera, a partir del currículo, la descripción de un empleo adecuado para la persona, como recomendación explicable \citep{ZhengGIRL2023}.

La revisión bibliométrica de \citet{RojasGaleano2022} sitúa estos desarrollos en un área en crecimiento, concentrada en el análisis de currículos, la extracción de habilidades y el \textit{ranking} con representaciones semánticas. Los sistemas revisados tienen en común que se entrenan y se evalúan con datos de interacción propietarios. En la búsqueda bibliográfica de este trabajo, los que combinan RAG y emparejamiento de empleo resultaron escasos y recientes.

\textbf{Relación con JobAI.} JobAI no entrena representaciones propias: usa un modelo de \textit{embeddings} general y un pipeline RAG con componentes disponibles. La pregunta que este trabajo responde es complementaria a la de estos sistemas: qué se obtiene, en este dominio, con componentes de uso general y sin datos de interacción. La extracción de restricciones de la consulta con un modelo de lenguaje sigue la misma idea que \citet{Liu2025}, a una escala mucho menor.

\section{Retrieval-Augmented Generation}
\label{sec:estado-arte-rag}

RAG fue propuesto por \citet{Lewis2020} como una arquitectura que combina un recuperador denso con un modelo generativo de secuencia a secuencia, entrenados de forma conjunta para responder a partir de pasajes recuperados de una memoria no paramétrica, en su caso Wikipedia. El diseño superó a los modelos puramente paramétricos en tareas intensivas en conocimiento, y la separación entre el modelo y el índice permite actualizar el conocimiento del sistema sin reentrenar el modelo.

La revisión de \citet{Gao2023} organiza la evolución posterior en tres niveles. El RAG \emph{ingenuo} indexa, recupera una vez y genera. El RAG \emph{avanzado} agrega optimizaciones antes de la recuperación, como la reformulación de la consulta, y después de ella, como el reordenamiento y la compresión del contexto. El RAG \emph{modular} descompone el pipeline en componentes intercambiables con flujos condicionales, como el enrutamiento entre estrategias de recuperación. La misma revisión trata, entre los recuperadores, la recuperación híbrida que combina representaciones dispersas y densas. Señala como retos abiertos la calidad de la recuperación, porque fragmentos irrelevantes degradan la respuesta aun con un buen generador, y la persistencia de alucinaciones cuando el contexto es insuficiente.

\textbf{Relación con JobAI.} Antes de este trabajo, la recuperación de JobAI era un RAG ingenuo: una sola búsqueda densa sin reordenamiento. Su enrutamiento entre una ruta semántica y una de agregación SQL ya correspondía a un rasgo del RAG modular (Sección~\ref{sec:estado-arte-hibrido}). La Iteración~2 incorpora componentes del RAG avanzado, y la evaluación mide el aporte de cada uno.

\section{Componentes del RAG avanzado}
\label{sec:estado-arte-componentes}

\subsection{Filtros estructurados y búsqueda aproximada con filtros}

Muchas consultas de empleo combinan una necesidad semántica con condiciones exactas: un cargo descrito con palabras libres, con un salario mínimo, en modalidad remota o con un tipo de contrato. Aplicar esas condiciones sobre los $k$ vecinos más cercanos ya recuperados (\emph{post-filtrado}) puede dejar la lista incompleta cuando el filtro es selectivo; aplicarlas antes, sobre toda la colección (\emph{pre-filtrado}), obliga a una búsqueda exacta sobre el subconjunto. Los índices de grafo como HNSW no se diseñaron para este caso. \citet{Gollapudi2023} proponen índices de grafo que incorporan las etiquetas de filtro en su construcción y en su búsqueda, y \citet{Patel2024} proponen ACORN, una extensión de HNSW que admite predicados arbitrarios recorriendo el subgrafo de los puntos que los cumplen. \citet{Sanca2024} estudian las rutas de acceso de la búsqueda mixta vector-relacional y muestran que el plan más eficiente depende de la selectividad del filtro y del volumen de datos.

\textbf{Relación con JobAI.} Las restricciones de la consulta se extraen con el clasificador y se aplican durante la búsqueda, tanto en la densa como en la léxica (Capítulo~\ref{cap:artefacto}). El \textit{benchmark} de almacenes mide cómo responde cada motor a filtros de distinta selectividad.

\subsection{Búsqueda híbrida y fusión de rankings}

La búsqueda híbrida combina una lista léxica y una densa con la intuición de que modelan la relevancia de formas complementarias: la léxica captura coincidencias exactas de términos poco frecuentes, como nombres de tecnologías, y la densa captura equivalencias sin términos comunes. La fusión por rangos recíprocos (RRF) de \citet{Cormack2009} suma, para cada documento, $1/(k + r)$ por cada lista en la que aparece en la posición $r$; no requiere calibrar los puntajes de los dos sistemas ni datos de entrenamiento, y sus autores mostraron que supera a otros métodos de fusión. \citet{Bruch2023}, en cambio, encontraron que RRF es sensible a su parámetro $k$ y que una combinación convexa de los puntajes normalizados la supera dentro y fuera del dominio, con pocos ejemplos para ajustar su único peso. En el emparejamiento de candidatos y ofertas, \citet{Yashwant2026} usan RRF para combinar la similitud vectorial con la búsqueda de texto completo, y comparan modelos de \textit{embeddings} con el puntaje del propio sistema de emparejamiento de su empresa.

\textbf{Relación con JobAI.} El pipeline fusiona la búsqueda densa y BM25 con RRF y $k = 60$, sin ajustar. Los hallazgos de \citet{Bruch2023} son una referencia para interpretar el aporte de este componente.

\subsection{Reordenamiento con \textit{cross-encoders}}

Los modelos densos codifican la consulta y el documento por separado y los comparan con una similitud entre vectores; un \textit{cross-encoder} los lee juntos y estima su relevancia, con más detalle y mucho mayor costo. Por eso se aplica solo a los primeros candidatos de una etapa barata. \citet{Nogueira2019} mostraron que un reordenador basado en BERT superaba en un 27\,\% relativo en MRR@10 al mejor sistema previo de MS MARCO. En BEIR, los modelos de reordenamiento obtuvieron los mejores resultados promedio, a un alto costo computacional \citep{Thakur2021}. Hoy existen reordenadores multilingües de tamaño moderado, como bge-reranker-v2-m3, basado en BGE-M3 \citep{Chen2024M3}, y la serie Qwen3-Reranker, construida sobre modelos de lenguaje de la familia Qwen3 \citep{Zhang2025Qwen3}.

\textbf{Relación con JobAI.} El pipeline reordena las 20 primeras candidatas con Qwen3-Reranker-0.6B, elegido frente a bge-reranker-v2-m3 en un conjunto de desarrollo (Capítulo~\ref{cap:artefacto}). La evaluación mide su aporte y su costo en latencia.

\subsection{Modelos de \textit{embeddings} multilingües}

La familia E5 \citep{Wang2022E5} se entrena de forma contrastiva con pares de textos débilmente supervisados; según sus autores, fue el primer modelo que superó a BM25 en BEIR sin datos etiquetados. Su versión multilingüe \citep{Wang2024mE5} sigue la misma receta con mil millones de pares de textos en varios idiomas, en tamaños pequeño, base y grande. BGE-M3 \citep{Chen2024M3} cubre más de 100 idiomas y combina recuperación densa, dispersa y de múltiples vectores en un solo modelo.

\textbf{Relación con JobAI.} Las ofertas del corpus están en inglés, pero los usuarios pueden escribir en español. La Iteración~2 reemplazó e5-base-v2, entrenado en inglés, por multilingual-e5-base, de la misma dimensión, y la evaluación compara ambos sobre las mismas consultas.

\section{Bases de datos vectoriales}
\label{sec:estado-arte-vectordb}

La recuperación densa requiere resolver la búsqueda de vecinos más cercanos sobre muchos vectores. \citet{Malkov2020} proponen HNSW, un grafo de proximidad organizado en capas jerárquicas con búsqueda de complejidad aproximadamente logarítmica, que hoy usan la mayoría de los motores vectoriales. \citet{Johnson2019}, con FAISS, aceleran en GPU la búsqueda exhaustiva y la aproximada mediante cuantización de producto. \citet{Aumuller2020}, con ANN-Benchmarks, proponen una metodología estandarizada para comparar estos algoritmos mediante curvas de compromiso entre \textit{recall} y velocidad.

Los algoritmos se ofrecen de dos formas. La primera es una extensión de un motor relacional, como pgvector sobre PostgreSQL, que agrega un tipo de dato vectorial e índices ANN. La segunda es un motor dedicado. \citet{Wang2021} describen con Milvus el diseño de esta segunda familia: separa almacenamiento y cómputo, admite varios tipos de índice y está pensada para cargas heterogéneas a gran escala. La revisión de \citet{Pan2024} organiza más de veinte de estos sistemas según su indexación, su soporte de búsqueda con filtros y su escalamiento, y la de \citet{Taipalus2024} señala como problemas abiertos la consistencia transaccional y la actualización incremental de los índices.

Esta literatura muestra que los motores dedicados están diseñados para escalas y patrones de acceso en los que una extensión relacional puede quedarse corta, en particular la búsqueda con filtros \citep{Sanca2024}. No establece, sin embargo, a partir de qué escala ni con qué selectividad la diferencia importa en la práctica, y una extensión relacional conserva ventajas operativas: transacciones, integridad referencial y una sola base de datos que administrar.

\textbf{Relación con JobAI.} JobAI almacenaba sus vectores en pgvector. La revisión del artefacto mostró que su lentitud se debía a una consulta que no usaba el índice y no al motor (Capítulo~\ref{cap:artefacto}). Por eso este trabajo no da por justificada la migración a Qdrant, sino que la evalúa: ambos motores quedan detrás de una interfaz común y se comparan con la metodología de \citet{Aumuller2020}, sobre datos reales y sintéticos, con y sin filtros, para establecer a partir de qué escala y con qué selectividad la elección del motor importa (Sección~\ref{sec:met-preguntas}).

\section{Clasificación de consultas y text-to-SQL}
\label{sec:estado-arte-hibrido}

No todas las preguntas sobre un catálogo de ofertas son de recuperación. Una pregunta como \textit{\enquote{¿cuántas ofertas remotas de Python hay?}} tiene una respuesta exacta que da una agregación SQL; resolverla con una búsqueda vectorial, que devuelve los fragmentos más similares y no un conteo, produce una respuesta aproximada. \citet{Biswal2024} formalizan esta observación con TAG (\textit{Table-Augmented Generation}): el text-to-SQL puro y el RAG puro cubren cada uno solo una parte de las preguntas que un usuario hace sobre una base de datos, y ambos caminos deben combinarse.

La dificultad de generar SQL correcto con un modelo de lenguaje sin restricciones queda documentada por \citet{Li2023} con BIRD, un banco de 12.751 pares de pregunta y consulta sobre 95 bases de datos reales de 37 dominios. Incluso modelos muy capaces quedan muy por debajo del desempeño humano cuando deben generar SQL ejecutable sobre esquemas reales, con sus inconsistencias y su volumen. Esa brecha motiva clasificar la intención de la consulta antes de decidir si se consulta la base de datos, en lugar de delegar todo a un modelo de lenguaje.

\textbf{Relación con JobAI.} El clasificador de JobAI combina patrones léxicos con un clasificador basado en LLM que elige entre recuperación y agregación y, en el segundo caso, una plantilla de un catálogo cerrado de consultas SQL parametrizadas (Capítulo~\ref{cap:artefacto}). El modelo nunca genera SQL libre, lo que evita el riesgo que documenta \citet{Li2023}, a cambio de cubrir solo las preguntas previstas en el catálogo. La evaluación mide la exactitud de ese enrutamiento.

\section{Evaluación de sistemas RAG}
\label{sec:estado-arte-evaluacion}

Evaluar un sistema RAG exige separar la recuperación de la generación. Para la recuperación, la práctica de las colecciones de prueba usa juicios de relevancia sobre un agrupamiento de los resultados de varios sistemas \citep{Sanderson2010}. Juzgar a mano todos los pares es costoso, y la literatura reciente estudia el uso de modelos de lenguaje como jueces. \citet{Thomas2024} mostraron que un LLM con un prompt cuidadosamente diseñado puede predecir las preferencias de evaluadores humanos con una precisión comparable a la de estos, y \citet{Upadhyay2024} publicaron UMBRELA, una reproducción abierta de ese procedimiento. \citet{Faggioli2023} advierten, sin embargo, sobre sus riesgos, entre ellos los sesgos sistemáticos y la circularidad cuando el juez y el sistema evaluado comparten modelo, y proponen una colaboración entre evaluadores humanos y modelos en lugar de un reemplazo completo.

Para la generación, \citet{Es2024} proponen RAGAS, un conjunto de métricas sin respuesta de referencia. La \emph{fidelidad} descompone la respuesta en afirmaciones y verifica cada una contra el contexto recuperado; la \emph{relevancia de la respuesta} mide si la respuesta atiende a la pregunta formulada.

La evaluación de un sistema RAG también debe considerar cómo falla. \citet{Barnett2024} reportan la experiencia de construir tres sistemas RAG en dominios distintos e identifican siete puntos de falla: la información no está en el corpus, el documento relevante no queda entre los primeros, no llega al contexto, el modelo no la extrae, la respuesta tiene un formato equivocado, una especificidad incorrecta o está incompleta. Concluyen que la validación de un sistema RAG solo es posible durante su operación y que su robustez se construye de forma evolutiva.

\textbf{Relación con JobAI.} La evaluación usa un modelo de lenguaje como juez, con el procedimiento de UMBRELA adaptado al dominio y validado contra juicios humanos, y mide la fidelidad y la relevancia de las respuestas con RAGAS. Además de medir la respuesta final, mide salidas intermedias del pipeline. Este trabajo sostiene, a partir de las fallas encontradas en JobAI, que un protocolo de evaluación por componentes revela antes de la operación una parte de las fallas que \citet{Barnett2024} solo observaron en ella.

\section{Design Science Research}
\label{sec:estado-arte-dsr}

Design Science Research (DSR) se distingue de la investigación conductual en sistemas de información en que no busca explicar o predecir un fenómeno, sino diseñar y evaluar con rigor un artefacto que resuelve un problema relevante. \citet{Hevner2004} lo formalizan en siete lineamientos, entre ellos que el artefacto es una contribución identificable, que su utilidad debe demostrarse con métodos de evaluación rigurosos y que la investigación debe comunicarse a audiencias técnicas y de gestión. Proponen además dos ciclos complementarios: el de \emph{relevancia}, que conecta el artefacto con las necesidades de su entorno, y el de \emph{rigor}, que lo conecta con la base de conocimiento existente. El volumen editado por \citet{HevnerChatterjee2010} consolida la teoría y la práctica de DSR con un tercer ciclo, el de \emph{diseño}, que construye y evalúa el artefacto de forma iterativa.

\citet{Peffers2007} operacionalizan estos principios en seis actividades: identificación del problema, definición de objetivos, diseño y desarrollo, demostración, evaluación y comunicación. Reconocen distintos puntos de entrada al ciclo, entre ellos uno centrado en el diseño y desarrollo, apropiado cuando la investigación parte de un artefacto existente. \citet{Venable2016} proponen FEDS, un marco para diseñar la evaluación en DSR que distingue la evaluación \emph{formativa}, que orienta el diseño mientras se construye, de la \emph{sumativa}, que juzga el resultado.

\textbf{Relación con JobAI.} Este trabajo adopta el punto de entrada centrado en el diseño y desarrollo, porque JobAI ya existía. Combina una evaluación formativa, cuyos hallazgos orientaron las dos iteraciones de diseño, con una evaluación sumativa de la configuración final (Capítulo~\ref{cap:metodologia}).

\section{Síntesis}

La literatura revisada permite tres afirmaciones. Primero, la búsqueda léxica tiene una limitación estructural en el reclutamiento, pero BM25 es una línea base robusta que un sistema denso no siempre supera. Segundo, cada componente del RAG avanzado tiene respaldo propio (los filtros en la búsqueda aproximada, la fusión híbrida, el reordenamiento y los modelos multilingües), pero su aporte se ha medido sobre todo en colecciones de propósito general, y para algunos, como la fusión híbrida, la evidencia es mixta. Tercero, las grandes plataformas de empleo usan representaciones semánticas aprendidas de sus propios datos, de modo que sus resultados no indican qué obtiene un agregador pequeño con componentes de uso general.

Falta, por tanto, una evaluación reproducible que mida, en la búsqueda de empleo, cuánto aporta cada componente frente a BM25 y cuánto cuesta, y que documente qué fallas de un sistema RAG se hacen visibles con una evaluación sistemática. El Capítulo~\ref{cap:metodologia} traduce este vacío en un protocolo de evaluación.
